% GENERATED FILE. Edit canonical lessons, then run npm run generate:books.
% canonical-source-hash: fnv1a64:f90e5291c008c8cb
% canonical-lessons: JA-C99-saifu, JA-C99-tsue, JA-C99-fukuro, JA-C99-tawashi, JA-C99-chiritori

\chapter{Wallet, Walking stick, Bag, Scrubbing brush, Dustpan}
\label{ch:ja-wallet-walking-stick-bag-scrubbing-brush-dustpan}
\hlchaptermodality{99}

\begin{chapteropening}
\textbf{By the end of this chapter:} \emph{I can say a wallet, a walking stick, a bag, a scrubbing brush and a dustpan.}

Complete the last lesson of chapter 99: I can say a wallet, a walking stick, a bag, a scrubbing brush and a dustpan.
\end{chapteropening}

\section[\texorpdfstring{\ja{さいふ} (saifu)}{さいふ (saifu)}]{\ja{さいふ} (saifu) — a wallet}

\label{lesson:JA-C99-saifu}



\begin{quote}
Before the new one: say the Japanese for the fault, then the Japanese for a reason.
\end{quote}

\subsection*{You'll want to know: \ja{さいふ}}
\textbf{\ja{さいふ}} — \emph{saifu} — \textquotedblleft{}a wallet\textquotedblright{}.

\begin{grammarlens}[title={What you've built}]
One more word for this chapter.
\end{grammarlens}

\subsection*{Your turn}
\begin{itemize}
\raggedright
  \item \emph{Say it:} \emph{saifu}
  \item \emph{Say it:} \emph{saifu}, once more
  \item \emph{Say it:} say \emph{saifu}
  \item \emph{Recall:} say the Japanese for the fault, then the Japanese for a reason, then say \emph{saifu} again
\end{itemize}

\begin{tcolorbox}[breakable,colback=teal!4,colframe=teal!35!black,title={Before you move on}]
What does \ja{さいふ} mean? (\textquotedblleft{}A wallet\textquotedblright{}.) Say it once more.
\end{tcolorbox}
\section[\texorpdfstring{\ja{つえ} (tsue)}{つえ (tsue)}]{\ja{つえ} (tsue) — a walking stick}

\label{lesson:JA-C99-tsue}



\begin{quote}
Before the new one: say the Japanese for a reason, then the Japanese for a wallet.
\end{quote}

\subsection*{You'll want to know: \ja{つえ}}
\textbf{\ja{つえ}} — \emph{tsue} — \textquotedblleft{}a walking stick\textquotedblright{}.

\begin{grammarlens}[title={What you've built}]
One more word for this chapter.
\end{grammarlens}

\subsection*{Your turn}
\begin{itemize}
\raggedright
  \item \emph{Say it:} \emph{tsue}
  \item \emph{Say it:} \emph{tsue}, once more
  \item \emph{Say it:} say \emph{tsue}
  \item \emph{Recall:} say the Japanese for a reason, then the Japanese for a wallet, then say \emph{tsue} again
\end{itemize}

\begin{tcolorbox}[breakable,colback=teal!4,colframe=teal!35!black,title={Before you move on}]
What does \ja{つえ} mean? (\textquotedblleft{}A walking stick\textquotedblright{}.) Say it once more.
\end{tcolorbox}
\section[\texorpdfstring{\ja{ふくろ} (fukuro)}{ふくろ (fukuro)}]{\ja{ふくろ} (fukuro) — a bag, a sack}

\label{lesson:JA-C99-fukuro}



\begin{quote}
Before the new one: say the Japanese for a wallet, then the Japanese for a walking stick.
\end{quote}

\subsection*{You'll want to know: \ja{ふくろ}}
\textbf{\ja{ふくろ}} — \emph{fukuro} — \textquotedblleft{}a bag, a sack\textquotedblright{}.

\begin{grammarlens}[title={What you've built}]
One more word for this chapter.
\end{grammarlens}

\subsection*{Your turn}
\begin{itemize}
\raggedright
  \item \emph{Say it:} \emph{fukuro}
  \item \emph{Say it:} \emph{fukuro}, once more
  \item \emph{Say it:} say \emph{fukuro}
  \item \emph{Recall:} say the Japanese for a wallet, then the Japanese for a walking stick, then say \emph{fukuro} again
\end{itemize}

\begin{tcolorbox}[breakable,colback=teal!4,colframe=teal!35!black,title={Before you move on}]
What does \ja{ふくろ} mean? (\textquotedblleft{}A bag, a sack\textquotedblright{}.) Say it once more.
\end{tcolorbox}
\section[\texorpdfstring{\ja{たわし} (tawashi)}{たわし (tawashi)}]{\ja{たわし} (tawashi) — a scrubbing brush}

\label{lesson:JA-C99-tawashi}



\begin{quote}
Before the new one: say the Japanese for a walking stick, then the Japanese for a bag.
\end{quote}

\subsection*{You'll want to know: \ja{たわし}}
\textbf{\ja{たわし}} — \emph{tawashi} — \textquotedblleft{}a scrubbing brush\textquotedblright{}.

\begin{grammarlens}[title={What you've built}]
One more word for this chapter.
\end{grammarlens}

\subsection*{Your turn}
\begin{itemize}
\raggedright
  \item \emph{Say it:} \emph{tawashi}
  \item \emph{Say it:} \emph{tawashi}, once more
  \item \emph{Say it:} say \emph{tawashi}
  \item \emph{Recall:} say the Japanese for a walking stick, then the Japanese for a bag, then say \emph{tawashi} again
\end{itemize}

\begin{tcolorbox}[breakable,colback=teal!4,colframe=teal!35!black,title={Before you move on}]
What does \ja{たわし} mean? (\textquotedblleft{}A scrubbing brush\textquotedblright{}.) Say it once more.
\end{tcolorbox}
\section[\texorpdfstring{\ja{ちりとり} (chiritori)}{ちりとり (chiritori)}]{\ja{ちりとり} (chiritori) — a dustpan}

\label{lesson:JA-C99-chiritori}



\begin{quote}
Before the new one: say the Japanese for a bag, then the Japanese for a scrubbing brush.
\end{quote}

\subsection*{You'll want to know: \ja{ちりとり}}
\textbf{\ja{ちりとり}} — \emph{chiritori} — \textquotedblleft{}a dustpan\textquotedblright{}.

\begin{grammarlens}[title={What you've built}]
One more word for this chapter.
\end{grammarlens}

\subsection*{Your turn}
\begin{itemize}
\raggedright
  \item \emph{Say it:} \emph{chiritori}
  \item \emph{Say it:} \emph{chiritori}, once more
  \item \emph{Say it:} say \emph{chiritori}
  \item \emph{Recall:} say the Japanese for a bag, then the Japanese for a scrubbing brush, then say \emph{chiritori} again
\end{itemize}

\begin{tcolorbox}[breakable,colback=teal!4,colframe=teal!35!black,title={Before you move on}]
What does \ja{ちりとり} mean? (\textquotedblleft{}A dustpan\textquotedblright{}.) Say it once more.
\end{tcolorbox}
